\chapter{A full quantitative counting argument from Hales--Jewett}
\label{qr:chapter}

The quantitative colouring input to Section 5 can be proved from the
Hales--Jewett theorem already in mathlib.  This chapter gives the
single-equation argument, including the finite Deuber construction used
to obtain the right number of independent parameters.
The starting point is the proof strategy of Frankl, Graham and R\"odl,
\emph{Quantitative Theorems for Regular Systems of Equations}, Journal of
Combinatorial Theory A 47 (1988), 246--261, Theorem 1.  Their Deuber variant
is attributed to an unpublished manuscript.  We instead expand the
Hales--Jewett row-focusing argument attributed to Leeb and Pr\"omel, as
presented by Joel Moreira in his 2014 exposition
\emph{A proof of Deuber's theorem using Hales--Jewett's theorem}.
Keeping track of finitely many disjoint-support embeddings replaces the
original residue encoding.  A separated generator box makes positivity
and containment in $[N]$ explicit.

The only Ramsey theorem invoked without a new proof here is the existing
mathlib theorem
\texttt{Combinatorics.Line.exists\_mono\_in\_high\_dimension}.
It states that every colouring of words over a fixed finite alphabet, in
a sufficiently high finite dimension, has a monochromatic combinatorial
line.  A line has a nonempty set of variable positions, all of which are
replaced by the same letter; the remaining positions have fixed letters.
Reindexing its finite coordinate type by $\{0,\ldots,h-1\}$ gives the
numerical form used below.  Taking any constant colouring shows that
the resulting $h$ is positive, because a line has a variable position.
The results in this chapter are detailed proof targets, not declarations
already checked in Lean.

\section{Finite patterns and their embeddings}

\begin{definition}[Rows of a Deuber pattern]
\label{qr:deuber-pattern}
For integers $m\geq0$, $p,c\geq1$ and $z=(z_0,\ldots,z_m)\in\mathbb Z^{m+1}$,
write $F_p=\{-p,\ldots,p\}$ and put
\[
 L(j,p,c;z)=\left\{cz_j+\sum_{i<j}\lambda_i z_i:
                                  \lambda_i\in F_p\right\},\quad
 D(m,p,c;z)=\bigcup_{j=0}^mL(j,p,c;z).
\]
In particular $L(0,p,c;z)=\{cz_0\}$ is a singleton.
We will only colour numerical patterns contained in the positive integers.
This condition implies $z_j>0$ for every $j$, because $cz_j$ belongs to
the pattern and $c>0$.
A colouring is $k$-regular, for $0\leq k\leq m$, if each row with
$m-k<j\leq m$ is monochromatic.  The colours of different rows may differ.
Every colouring is $0$-regular; an $m$-regular colouring has all its rows
monochromatic, since row zero is a singleton.
\end{definition}

\begin{definition}[Disjoint-support generator embeddings]
\label{qr:block-embedding}
An integer matrix $T$ with $q$ rows and $M$ columns is a block embedding if
each row has a nonempty support and the supports of distinct rows are
disjoint.  For a seed vector $Y\in\mathbb Z^M$, its new generators are
$y=TY$.  A finite family of such matrices can depend on the pattern
parameters and the number of colours, but is chosen before the seed
values and their colouring.
\end{definition}

\begin{lemma}[Composition preserves the block property]
\label{qr:block-composition}
\uses{qr:block-embedding}
If $T$ and $U$ are block embeddings and their matrix product is defined,
then $TU$ is a block embedding.  Selecting distinct rows of a block
embedding also gives a block embedding.
\end{lemma}
\begin{proof}
For a fixed source coordinate $b$, at most one row $a$ of $U$ has
$U_{ab}\ne0$, because its row supports are disjoint.
Thus $(TU)_{ib}=\sum_aT_{ia}U_{ab}$ has at most one nonzero summand.
It is nonzero exactly when $b$ belongs to the support of some row $a$ of
$U$ for which $T_{ia}\ne0$; there can be no cancellation.
Consequently the support of row $i$ of $TU$ is the union of those row
supports of $U$ selected by the support of row $i$ of $T$.
This union is nonempty, since both matrices have nonempty row supports.
For different $i$ these unions use disjoint sets of rows of $U$, and
are therefore disjoint.  Row selection preserves the property directly.
\end{proof}

\begin{lemma}[One additional monochromatic row]
\label{qr:one-step}
\uses{qr:deuber-pattern,qr:block-embedding}
Fix $m\geq1$, $p,c,r\geq1$ and $0\leq k<m$.
There are integers $M\geq m$, $P,C\geq1$ and a finite family of block
embeddings from $M+1$ generators to $m+1$ generators with the following
property.  If $D(M,P,C;S)$ is positive and its $r$-colouring is $k$-regular,
some matrix in this family produces generators $s$ for which
$D(m,p,c;s)$ is a subset of $D(M,P,C;S)$ and is $(k+1)$-regular.
\end{lemma}
\begin{proof}
Put $d=m-k\geq1$.  Apply Hales--Jewett to the alphabet $F_p^d$ and $r$
colours and choose a positive word length $h$ with its conclusion.
Set
\[
 b=hd,\qquad M=b+k,\qquad
 P=\max(p^2,pc,c^2),\qquad C=c^2.
\]
Since $h\geq1$, $M\geq d+k=m$.
For a word $a=(a_0,\ldots,a_{h-1})$, with each
$a_i=(a_{i,0},\ldots,a_{i,d-1})\in F_p^d$, define its colour to be the
colour of
\[
 c\sum_{i<h}\sum_{t<d}a_{i,t}S_{id+t}+c^2S_b.                 \tag{1}
\]
This integer lies in row $b$ of the source pattern, since its earlier
coefficients have absolute value at most $cp\leq P$ and its leading
coefficient is $C$.  The induced word colouring is therefore well-defined.

Choose a monochromatic line.  Let $A\subset\{0,\ldots,h-1\}$ be its
nonempty set of variable positions and let $B$ be the complement.
For $i\in B$ write $w_i\in F_p^d$ for its fixed letter.
Define the new generators by
\[
 s_t=c\sum_{i\in A}S_{id+t}\quad(0\leq t<d),
\]
\[
 s_d=\sum_{i\in B}\sum_{t<d}w_{i,t}S_{id+t}+cS_b,
 \qquad
 s_t=cS_{b+t-d}\quad(d<t\leq m).                            \tag{2}
\]
These are integer linear combinations of the source generators.
For each early row $t<d$, its support is
$\{id+t:i\in A\}$ and is nonempty.  These early supports are mutually
disjoint.  The middle generator uses only the complementary $B$-blocks
and the new coordinate $b$, where its coefficient is $c\ne0$.
The last $k$ generators use respectively the single coordinates
$b+1,\ldots,b+k$.  All these supports are disjoint.
Thus (2) defines a block embedding.  There are only finitely many choices
of the nonempty set $A$ and fixed letters $w_i$, so all resulting
embeddings form a fixed finite family.

We verify every row inclusion, not only the row made monochromatic.
For $d<j\leq m$, a point of row $j$ of the new pattern has leading term
$cs_j=c^2S_{b+j-d}$.
Every generator $s_t$ with $t<j$ uses source indices strictly below
$b+j-d$.  Because the generator supports are disjoint, their coefficients
do not add at any coordinate; each earlier coefficient has modulus at
most $p\max(p,c)\leq P$.
Hence
\[
 L(j,p,c;s)\subset L(b+j-d,P,C;S).                           \tag{3}
\]
These are the last $k$ rows of the source pattern, so they are
monochromatic by its $k$-regularity.

For row $j=d$, expand using (2): its points are
\[
 c\sum_{t<d}\sum_{i\in A}\lambda_tS_{id+t}
 +c\sum_{i\in B}\sum_{t<d}w_{i,t}S_{id+t}+c^2S_b,
 \qquad(\lambda_0,\ldots,\lambda_{d-1})\in F_p^d.
\]
These are exactly the integers (1) associated with the chosen
combinatorial line.  They are monochromatic and lie in source row $b$.

Finally let $0\leq j<d$, and let $a=\max A$, which exists because
$A$ is finite and nonempty.  A point of new row $j$ is
\[
 c\sum_{t<j}\lambda_t\sum_{i\in A}S_{id+t}
       +c^2\sum_{i\in A}S_{id+j}.
\]
Its largest source index is $ad+j$, with coefficient $c^2=C$.
All other nonzero coefficients have earlier indices and modulus at most
$\max(cp,c^2)\leq P$.  It therefore belongs to source row $ad+j$.
This proves containment for every remaining row.
Every new row is consequently positive, as it is contained in the
positive source pattern; in particular the new generators are positive.
The last $k$ new rows are monochromatic by (3), and the next row $d$ is
monochromatic by Hales--Jewett.  The new colouring is $(k+1)$-regular.
\end{proof}

\begin{theorem}[Finite Deuber extraction with finitely many block embeddings]
\label{qr:finite-deuber}
\uses{qr:one-step,qr:block-composition}
For every number of generators $q\geq1$ and integers $p,c,r\geq1$, there
are $M\geq q$, $P,C\geq1$, and a finite nonempty family
$\mathcal F$ of $q$-by-$M$ block embeddings such that every $r$-colouring
of a positive pattern $D(M-1,P,C;Y)$ contains a monochromatic pattern
\[
 D(q-1,p,c;TY)\subset D(M-1,P,C;Y)
\]
for some $T\in\mathcal F$.
\end{theorem}
\begin{proof}
If $q=1$, take $M=1$, $P=p$, $C=c$, and the identity embedding.
The target is its single row $\{cY_0\}$, which is monochromatic.
Assume $q\geq2$ and set $n=r(q-1)\geq1$.
Start with $(m_0,p_0,c_0)=(n,p,c)$.
For $j=1,\ldots,n$, apply Lemma~\ref{qr:one-step} with target parameters
$(m_{j-1},p_{j-1},c_{j-1})$ and $k=n-j$, and let its source parameters
be $(m_j,p_j,c_j)$.  This recursion is legitimate: $m_j\geq m_{j-1}\geq n$
and $0\leq n-j<n\leq m_{j-1}$.
All parameters and finite embedding families have now been chosen solely
from $q,p,c,r$.

Take a colouring of a positive $D(m_n,p_n,c_n;Y)$.
It is $0$-regular.  Use the step with $j=n$ to obtain a $1$-regular
$(m_{n-1},p_{n-1},c_{n-1})$ pattern inside it.
Use the step with $j=n-1$ to obtain a $2$-regular pattern inside that one.
Continuing for all $n$ steps produces an $n$-regular
$D(n,p,c;z)$ in the original source pattern.
Every one of its $n+1$ rows is monochromatic, including the singleton
row zero.  Since $n+1=r(q-1)+1$, the pigeonhole principle supplies $q$
distinct rows of the same colour, with indices
$0\leq j_0<\cdots<j_{q-1}\leq n$.
Put $y_i=z_{j_i}$.  For each target row $i$, a combination
$cy_i+\sum_{t<i}\lambda_t y_t$ belongs to source row $j_i$: set the
coefficients of every unselected earlier generator to zero.
Thus $D(q-1,p,c;y)$ is monochromatic and lies in $D(n,p,c;z)$.

Every generator change was a matrix from a fixed finite family of block
embeddings.  Composition and the final row selection preserve that property
by Lemma~\ref{qr:block-composition}.
There are finitely many compositions and finitely many row selections.
Let $\mathcal F$ be their set of resulting matrices, and put
$M=m_n+1$, $P=p_n$, $C=c_n$.
This family is nonempty, since each Hales--Jewett step has at least one
possible line and the indicated row selections exist for the constructed
colourings.  Each matrix in it has $q$ nonempty disjoint row supports, so
choosing one coordinate in each support injects $q$ rows into $M$ columns;
hence $M\geq q$.
The preceding construction proves the claimed extraction for every seed
and every colouring.
\end{proof}

\section{A full-dimensional first-entry kernel parametrization}

\begin{lemma}[A kernel matrix whose rows are Deuber coefficients]
\label{qr:kernel-matrix}
Let $a_1,\ldots,a_s\in\mathbb Z\setminus\{0\}$, and suppose
$\sum_{i\in I}a_i=0$ for some nonempty $I\subset\{1,\ldots,s\}$.
Then $s\geq2$, and for $q=s-1$ there are $p,c\geq1$ and an integer
matrix $B$ with $s$ rows and $q$ columns such that:
\begin{enumerate}
\item $a^{\mathsf T}B=0$;
\item the map $B:\mathbb Q^q\to\mathbb Q^s$ is injective;
\item each row of $B$ has last nonzero entry $c$, and every earlier
      entry has absolute value at most $p$.
\end{enumerate}
Consequently each coordinate of $By$ belongs to $D(q-1,p,c;y)$.
\end{lemma}
\begin{proof}
A singleton $I$ cannot have zero sum because every $a_i$ is nonzero.
Thus $s\geq2$ and $q\geq1$.
Work first over $\mathbb Q$ in the kernel
$V=\{v:\sum_i a_iv_i=0\}$.
For any pivot $i_0$, solving for $v_{i_0}$ gives a linear isomorphism
between $V$ and $\mathbb Q^{s-1}$; its dimension is therefore $q$.

The vector $u_0=1_I$ is a nonzero element of $V$.
If $I$ is all indices, retain this one vector as the beginning of a basis.
If $I$ is proper, choose $i_0\in I$ and define $u_1$ by setting it equal
to one outside $I$, zero on $I\setminus\{i_0\}$, and
\[
 (u_1)_{i_0}=-\frac{\sum_{i\notin I}a_i}{a_{i_0}}.
\]
Its weighted coordinate sum is zero, so $u_1\in V$.
The pair $u_0,u_1$ is linearly independent: any relation first evaluated
outside $I$ forces the coefficient of $u_1$ to vanish, and then a
coordinate inside $I$ forces the coefficient of $u_0$ to vanish.

Extend this ordered independent list to a basis of $V$.
In the matrix with those basis vectors as columns, each row has first
nonzero entry one: for an index in $I$ it is in column $u_0$;
for an index outside $I$ it is in column $u_1$.
Every index is covered by one of these alternatives.
Clear the finitely many rational denominators by multiplying the whole
matrix by one positive integer $c$.
The matrix becomes integral, is still an injective parametrization of
the kernel, and its row first entries become the common positive integer
$c$.  Reverse its column order to obtain $B$.
Now each row has last nonzero entry $c$.
Choose $p$ to be the maximum of one and all absolute values of its entries.
All earlier entries are then in $F_p$, and the entries after the indicated
last column are zero.  This is exactly the coefficient form defining a
row of $D(q-1,p,c;y)$, proving the last assertion as well.
\end{proof}

\section{Many independent seed generators and bounded multiplicity}

\begin{lemma}[A positive generator box of the correct size]
\label{qr:positive-box}
\uses{qr:deuber-pattern}
Let $M,P,C\geq1$, put
\[
 R=4P+2,\qquad H=(2C+1)R^M,
\]
and let $N\geq2H$ be an integer.  There is a set $\mathcal Y_N$ of at
least $(N/(2H))^M$ integer vectors $Y\in[N]^M$ such that
$D(M-1,P,C;Y)\subset[N]$ for every $Y\in\mathcal Y_N$.
\end{lemma}
\begin{proof}
Let $L=\lfloor N/H\rfloor\geq N/(2H)\geq1$.
Independently choose each generator $Y_j$ in the integer interval
$[LR^j,2LR^j]$, for $j=0,\ldots,M-1$.
Each interval has at least $L$ choices, so there are at least
$L^M\geq(N/(2H))^M$ vectors.
The lower endpoint is positive.  The upper endpoint is at most $N$, since
$2LR^{M-1}\leq2NR^{M-1}/H\leq N$.

For a point of pattern row $j$, the absolute contribution of all earlier
terms is at most
\[
 P\sum_{i<j}Y_i\leq2PL\sum_{i<j}R^i
   =2PL\frac{R^j-1}{R-1}<\frac12 LR^j,
\]
because $R-1=4P+1>4P$.  For $j=0$ the earlier sum is zero and the strict
bound is still true.
The leading term $CY_j$ is at least $CLR^j$, so the whole expression is
positive.  Its upper bound is
\[
 2CLR^j+\tfrac12LR^j
 <(2C+1)LR^j
 \leq(2C+1)\frac NH R^{M-1}=N/R\leq N.
\]
Every pattern entry is an integer between one and $N$, as required.
\end{proof}

\begin{lemma}[The multiplicity bound]
\label{qr:multiplicity}
\uses{qr:block-embedding}
Let $B:\mathbb Q^q\to\mathbb Q^s$ be injective, and let
$\mathcal F$ be a finite family of $q$-by-$M$ integer block embeddings.
For any fixed $x\in\mathbb Z^s$ and integer $N\geq1$, the number of
$Y\in[N]^M$ for which $BTY=x$ for some $T\in\mathcal F$ is at most
\[
 |\mathcal F|N^{M-q}.
\]
\end{lemma}
\begin{proof}
Fix $T$.  If it has any such preimages, injectivity of $B$ implies that
their intermediate values $TY$ are all the same vector $y$.
Choose a pivot coordinate $j_i$ in the nonempty support of row $i$ of
$T$.  Disjointness makes these $q$ coordinates distinct and makes the
coefficient at $j_i$ zero in every other row.
Once the remaining $M-q$ coordinates of $Y$ are prescribed, the equation
$\sum_jT_{ij}Y_j=y_i$ determines $Y_{j_i}$ uniquely, if it has any integer
solution at all: its coefficient $T_{i,j_i}$ is nonzero, and the other
pivot coordinates have coefficient zero in that equation.
Equivalently, two solutions agreeing off the pivots have
$T_{i,j_i}(Y_{j_i}-Y'_{j_i})=0$, and cancellation in $\mathbb Z$ gives
equality at the pivot as well.
Restriction to nonpivot coordinates is therefore injective on the fiber.
There are $N^{M-q}$ such restrictions.
Finally take the union over the finite family $\mathcal F$; its cardinality
is at most the sum of these bounds.  Duplicate preimages only make the
union smaller, so no uniqueness of $T$ is needed.
\end{proof}

\begin{theorem}[Quantitative Rado for one homogeneous equation]
\label{qr:many-monochromatic}
\uses{qr:finite-deuber,qr:kernel-matrix,qr:positive-box,qr:multiplicity}
Let $a_1,\ldots,a_s\in\mathbb Z\setminus\{0\}$ have a nonempty zero-sum
subcollection, and let $r\geq1$.  There are $N_0=N_0(a,r)$ and
$\delta_0=\delta_0(a,r)>0$ such that every $r$-colouring of $[N]$, for
$N\geq N_0$, has a colour class $C_j$ with
\[
 \#\{x\in C_j^s:\textstyle\sum_i a_i x_i=0\}
 \geq\delta_0N^{s-1}.
\]
In particular this proves the previously external input in
Lemma~\ref{app:fgr-dependency} from the existing Hales--Jewett theorem.
\end{theorem}
\begin{proof}
Fix $q=s-1$, and choose the integral injective kernel parametrization
$B$ and its parameters $p,c$ from Lemma~\ref{qr:kernel-matrix}.
Use Theorem~\ref{qr:finite-deuber} to fix $M,P,C$ and a finite nonempty
family $\mathcal F$ of block embeddings.
Set $F=|\mathcal F|\geq1$, $R=4P+2$, $H=(2C+1)R^M$, and $N_0=2H$.
All these choices depend only on the equation and $r$.

For each $N\geq N_0$ use Lemma~\ref{qr:positive-box} to obtain
$\mathcal Y_N\subset[N]^M$ with
$|\mathcal Y_N|\geq(N/(2H))^M$.
Restrict the given colouring of $[N]$ to the positive pattern associated
with each seed $Y\in\mathcal Y_N$.
The finite Deuber theorem supplies some $T\in\mathcal F$ for which
$D(q-1,p,c;TY)$ is monochromatic and is contained in $[N]$.
Set $x=BTY$.
Every coordinate of $x$ belongs to that target pattern, by the kernel
matrix construction, so $x$ is monochromatic and has all coordinates in
$[N]$.  Moreover $a^{\mathsf T}x=a^{\mathsf T}BTY=0$.
Choose one such output tuple for each seed; all sets involved are finite.

By Lemma~\ref{qr:multiplicity}, a fixed output tuple is selected for at
most $FN^{M-q}$ seeds.  Counting the fibers of this map gives at least
\[
 \frac{|\mathcal Y_N|}{FN^{M-q}}
 \geq \frac{N^q}{F(2H)^M}
\]
distinct monochromatic solution tuples in total.
Each belongs to a unique colour class, since $s\geq1$ and the classes
are disjoint.  One of the $r$ classes therefore contains at least a
proportion $1/r$ of these tuples.
Take
\[
 \delta_0=\frac1{rF(2H)^M}>0.
\]
Since $q=s-1$, this is exactly the claimed bound.
Repeated coordinates are allowed throughout; no extra distinctness
hypothesis was used or obtained.
\end{proof}

\section{What the construction makes explicit}

Every finite existence choice in this proof precedes the ambient length
$N$ and its colouring.  The kernel basis and denominator clearing determine
$B,p,c$.  A finite recursion using Hales--Jewett lengths determines
$M,P,C$ and the finite embedding family.  The elementary formulas above
then give $N_0$ and $\delta_0$.  Numerical optimization is unnecessary for
the Sidon application, whose constants may depend arbitrarily on the
coefficient vector and the number of colours.

The new proof does not invoke an arithmetic removal lemma, a new infinite
Ramsey theorem, or the unpublished Deuber statement in the 1988 paper.
Its additional formalization work consists of finite sets, finite words,
integer matrices, basis extension over $\mathbb Q$, and counting fibers.
The hardest combinatorial theorem, Hales--Jewett, is already checked in
mathlib.  It remains necessary to implement and check these reductions;
this natural-language chapter is not itself a machine-checked proof.
